\documentclass[11pt,a4paper]{article}
\usepackage[T1]{fontenc}
\usepackage[utf8]{inputenc}
\usepackage{lmodern}
\usepackage[margin=27mm,headheight=15pt]{geometry}
\usepackage{amsmath,amssymb,amsthm,mathtools}
\usepackage{booktabs,longtable,array}
\usepackage{xcolor,microtype,enumitem,fancyhdr,hyperref,listings}
\hypersetup{colorlinks=true,linkcolor=blue!45!black,urlcolor=blue!50!black,
 citecolor=blue!45!black,pdftitle={The planar isoperimetric inequality in iPRA},
 pdfauthor={Research note prepared for Vladimir Reshetnikov}}
\pagestyle{fancy}
\fancyhf{}
\fancyhead[L]{\small Isoperimetry in intuitionistic PRA}
\fancyhead[R]{\small Finitary proof and Lean blueprint}
\fancyfoot[C]{\thepage}
\setlength{\parskip}{5pt}
\setlength{\parindent}{0pt}
\setlist{itemsep=3pt,topsep=5pt}
\allowdisplaybreaks
\newtheorem{theorem}{Theorem}[section]
\newtheorem{lemma}[theorem]{Lemma}
\newtheorem{proposition}[theorem]{Proposition}
\newtheorem{corollary}[theorem]{Corollary}
\theoremstyle{definition}
\newtheorem{definition}[theorem]{Definition}
\newtheorem{remark}[theorem]{Remark}
\newcommand{\ipra}{\mathsf{iPRA}}
\newcommand{\N}{\mathbb N}
\newcommand{\Q}{\mathbb Q}
\newcommand{\C}{\mathbb C}
\newcommand{\eps}{\varepsilon}
\newcommand{\abs}[1]{\left|#1\right|}
\newcommand{\norm}[1]{\left\|#1\right\|}
\newcommand{\Area}{\mathcal A}
\newcommand{\len}{\mathcal L}
\newcommand{\sgn}{\operatorname{sgn}}
\newcommand{\sd}{\operatorname{sd}}
\newcommand{\Rea}{\operatorname{Re}}
\newcommand{\Ima}{\operatorname{Im}}
\newcommand{\dH}{d_{\mathrm H}}
\newcommand{\PR}{\mathsf{PR}}
\lstset{basicstyle=\ttfamily\small,breaklines=true,columns=fullflexible,
 keepspaces=true,frame=single,rulecolor=\color{black!20},
 backgroundcolor=\color{black!2}}

\begin{document}
\begin{titlepage}
\vspace*{12mm}
{\Large\bfseries The sharp planar isoperimetric inequality\\[3mm]
in intuitionistic primitive recursive arithmetic\par}
\vspace{8mm}
{\large A complete finitary human proof, effective equality analysis,\\
and a path to an internal Lean formalization\par}
\vspace{10mm}
Prepared for Vladimir Reshetnikov\\
14 November 2026
\vspace{10mm}
\begin{abstract}
We give an arithmetic proof of the classical planar inequality
\(L^2\ge4\pi A\). Its uniform finite core is a theorem about rational
polygons: if \(u_i\) are certified rational upper bounds on the side
lengths, \(S=\sum_i u_i\), and \(k\in\N\), then
\[
 4\pi_k\,|\Area(P)|\le S^2+2^{-k}.
\]
The numbers \(\pi_k\) are explicitly defined rationals satisfying
\(0\le\pi-\pi_k<2^{-k}\). The proof uses finite Fourier sums on
dyadic regular polygons, positive quadratic radicals, and a bounded
rational subdivision. All analytic limiting language is replaced by
precision estimates.

The finite theorem extends to polygonally presented rectifiable
boundaries. A lower-real formulation of perimeter avoids assuming that
the length of a computable curve is computable. An explicit Fourier
stability estimate proves that equality forces the boundary to be a
circle; its center and radius are extracted without compactness or a
choice of subsequence.

The formalization target is an actual derivation in ProveIt's embedded
intuitionistic PRA, not merely a Lean theorem about real numbers.
The included Lean files specify the finite data, numerical coding,
and internal target. They are an uncompiled blueprint, not a completed
machine-checked proof.
\end{abstract}
\vfill
\begin{tabular}{@{}p{0.23\textwidth}p{0.70\textwidth}@{}}
\toprule
Scope & The classical \emph{planar} inequality, its sharp constant,
and circle rigidity for the stated effective geometric presentations.\\
Human argument & Complete finite proof with explicit approximation
bounds and an induction audit for \(\ipra\).\\
Lean status & Concrete definitions and proof obligations; internal
derivation not yet constructed or compiled.\\
Other dimensions & Higher-dimensional, spherical, and metric-space
generalizations are outside this note.\\
\bottomrule
\end{tabular}
\end{titlepage}
\setcounter{tocdepth}{1}
\tableofcontents
\newpage

\section{What it means to prove this theorem in \texorpdfstring{\(\ipra\)}{iPRA}}

The usual planar theorem says that a rectifiable Jordan curve of
length \(L\), enclosing area \(A\), satisfies
\[
 L^2\ge4\pi A,
\]
with equality exactly for circles. This is the classical form discussed
by Osserman \cite{osserman}. The referenced encyclopedia page also
discusses many other inequalities; we address this planar theorem.

First-order arithmetic does not quantify over arbitrary subsets of
\(\mathbb R^2\). Therefore a claim of a proof in \(\ipra\)
must state how the geometry is represented. We first prove one uniform
first-order theorem about finite numerical codes. Curves are then
precision-indexed abbreviations for arithmetic statements about those
codes. No axiom asserting the existence of a completion of the rationals
is used.

\subsection{The uniform arithmetic theorem}
For a nonempty rational polygon
\(P=(v_0,\ldots,v_{m-1})\), put \(v_m=v_0\), \(e_i=v_{i+1}-v_i\), and
\[
 \Area(P)=\frac12\sum_{i<m}\det(v_i,v_{i+1}).
\]
Simplicity is not required: this is signed shoelace area.
An admissible upper-length list is a rational list
\(u=(u_0,\ldots,u_{m-1})\) satisfying
\[
 u_i\ge0,\qquad \norm{e_i}^2\le u_i^2.
 \tag{1.1}\label{eq:input}
\]
Thus \(u_i\) bounds a \emph{length}; its square bounds squared length.
Let \(S=\sum_{i<m}u_i\).

\begin{theorem}[Uniform finite isoperimetry]\label{thm:finite}
There is a fixed primitive recursive rational function \(k\mapsto\pi_k\),
defined in Section~\ref{sec:pi}, for which \(\ipra\) proves
\[
 \boxed{4\pi_k\,|\Area(P)|\le S^2+2^{-k}}
 \tag{1.2}\label{eq:finite}
\]
for every natural \(k\), every nonempty rational polygon code \(P\),
and every admissible upper-length code \(u\).
In the circle interpretation,
\[
 0\le\pi-\pi_k<2^{-k}.
 \tag{1.3}\label{eq:piaccuracy}
\]
\end{theorem}
After denominators have been cleared, the hypotheses and conclusion of
\eqref{eq:finite} are primitive recursive Boolean predicates of three
natural-number inputs. Its universal closure is an ordinary first-order
arithmetical sentence. Degenerate, repeated-vertex, clockwise, and
self-intersecting polygons are all included.

\subsection{Geometric interpretation and its precise scope}
A convenient effective boundary presentation consists of rational
polygonal approximations, explicit uniform or Hausdorff error bounds,
an area approximation modulus, and, when a Cauchy-real perimeter is used,
a perimeter approximation modulus. Such data turn the classical
inequality into a schema of arithmetic inequalities. We prove that
schema, including rigidity.

There is a more general formulation. A primitive recursive presentation
of a continuous closed curve, with a supplied uniform continuity modulus
and a variation bound, has a computable signed-area name. Its length may
be represented by the lower cut generated by the lengths of finite
inscribed polygons; no length convergence modulus is then required.
Section~\ref{sec:lower} proves the inequality in this representation.

For an \emph{effective Jordan domain}, ordinary enclosed area is supplied
by compatible simple polygonal domain approximations. Shoelace area is
then its area name. One may equivalently use a separately justified
effective Jordan/Green representation bridge. This note does not claim
that bare pointwise injectivity automatically supplies all such
computational data, nor that every classically specified curve has a
primitive recursive name. That assertion is neither the language nor
the content of the arithmetic theorem.

These distinctions matter. Computable simple rectifiable curves can
fail to have computable constant-speed parametrizations
\cite{gu-lutz-mayordomo}. The proof below never extracts a general
arclength parametrization from bare rectifiability.

\subsection{What is established, and what remains for Lean}
The human proof gives the finite identities, all mesh and precision
bounds, and the bounded induction invariants used to read the argument
in \(\ipra\). It does not provide a fully expanded natural-deduction
proof term. That syntactic task is the remaining Lean formalization.
In particular, the absence of classical choice from an external Lean
theorem would not by itself establish provability in \(\ipra\).

\section{Finite arithmetic and positive radicals}\label{sec:arithmetic}
\subsection{The arithmetic substrate}
Use intuitionistic first-order logic with equality, symbols for
primitive recursive functions, their defining equations, the successor
axioms, and quantifier-free induction. This is the ordinary
\(\ipra\) presentation implemented in \texttt{ConciseIPRA.PR}
\cite{proveit-pr}.

Integers are coded by signs and magnitudes, rationals by an integer
numerator and a positive denominator, and finite lists by natural
numbers. Rational addition, multiplication, comparison, and floor are
primitive recursive. Comparisons are decidable because they reduce to
integer comparisons after multiplication by positive denominators.
No decision procedure for arbitrary real numbers is being assumed.

A bounded conjunction
\(\bigwedge_{i<n}R(i)\), with primitive recursive Boolean \(R\), is
represented by the value of a Boolean fold. It is therefore a
quantifier-free assertion in the language with those primitive recursive
symbols. The same observation applies to bounds on every entry of a
finite table. Section~\ref{sec:audit} identifies the actual invariants;
the fact that a computation is primitive recursive alone does not prove
its correctness in \(\ipra\).

\subsection{Square-root intervals}
For rational \(x\ge0\) and \(t\in\N\), define
\[
 \sd(x,t)=2^{-t}\max\{j\le M:j^2\le2^{2t}x\},
 \tag{2.1}\label{eq:sd}
\]
where any explicit integer bound \(M>2^t\sqrt{x}\) may be used.
For \(0\le x\le1\), \(M=2^t\) suffices, allowing the endpoint.
For arbitrary \(x=a/b\ge0\), \(M=2^t(a+1)\) suffices.
The maximum is computed by bounded search, and
\[
 \sd(x,t)^2\le x<\bigl(\sd(x,t)+2^{-t}\bigr)^2.
 \tag{2.2}\label{eq:sdbracket}
\]
The proof is the loop invariant that the stored index is the greatest
accepted index seen so far, together with monotonicity of integer squares.

The reference Python uses integer square root to compute the same
maximum efficiently. The Lean blueprint uses a literal bounded maximum,
so its primitive recursive character is visible.

\subsection{Exact radical notation is finite algebra}
For readability, the finite proof uses positive square roots such as
\(\sqrt{(1+c)/2}\). They need not be treated as arbitrary real numbers.
They can be coded in a finite tower of positive quadratic radicals.
At one level an expression has the form \(a+b\sqrt d\), where \(a,b,d\)
belong to the preceding level and \(d\ge0\).

Addition and multiplication are finite coefficient-array operations.
Its sign is determined recursively. If \(d=0\), use the sign of \(a\).
Handle zero coefficients directly. If the nonzero coefficients \(a,b\)
have the same sign, their sum has that sign. If they have opposite signs,
\[
 \sgn(a+b\sqrt d)
   =\sgn(a)\,\sgn(a^2-b^2d).
 \tag{2.3}\label{eq:sign}
\]
Indeed, compare the two nonnegative magnitudes
\(|a|\) and \(|b|\sqrt d\) by comparing their squares.
Every recursive sign call in \eqref{eq:sign} drops one tower level.

This also handles redundant extensions: equality is determined by
sign zero, not by equality of coefficient arrays. One must not blindly
invert a formal element by its conjugate; the associated formal norm
can vanish in a redundant extension. The main Fourier identity avoids
algebraic division altogether, apart from rational scalar division.
For the rigidity argument, divisions are by quantities already proved
strictly positive, and may be represented as numerator/positive-denominator
pairs.

The algorithms are fixed recursions on finite arrays. Array lengths
double at each radical level; coefficient growth is bounded by iterating
the explicit integer bounds for addition and multiplication.
Nested loops can be unrolled to a bound obtained by primitive recursion.
There is no invocation of a universal evaluator for arbitrary primitive
recursive programs.

The next lemma proves the order laws behind this recursion, rather than
assuming that a terminating sign program automatically has them.

\subsection{A constructive ordered-extension lemma}
\begin{lemma}[Positive quadratic extension by finite certificates]
\label{lem:orderedextension}
Let \(K\) be a decidably linearly ordered commutative integral domain,
and let \(d\ge0\) belong to \(K\). Pairs \((a,b)\) with
\[
 (a,b)(c,e)=(ac+bed,\,ae+bc)
\]
admit the recursive sign rule \eqref{eq:sign}. Equality defined by sign
zero of the difference is a congruence. On the resulting classes it
defines a linearly ordered integral domain extending \(K\), and
\((0,1)\) is the nonnegative square root of \(d\).
All witnesses in the proof below are explicit finite expressions.
\end{lemma}
\begin{proof}
If \(d=0\), assign \((a,b)\) the sign of \(a\). The projection
\((a,b)\mapsto a\) proves every assertion. Hence suppose \(d>0\).
Let \(R\) be the pair ring; it is \emph{not} assumed to be a domain.
Define
\[
 T=\{(a,b):a\ge0,\ b\ge0,\ a+b>0\}.
\]
It contains \(1=(1,0)\) and \(r=(0,1)\), excludes zero, and is closed
under addition and multiplication. For multiplication, all coefficients
are nonnegative. Choose a positive coefficient in each factor; the
four possibilities give a positive contribution to \(ac+bed\) or
\(ae+bc\), using \(d>0\) in the \(be\) case.
These choices are decisions on two finite pairs.

For every \(x=(a,b)\), explicitly compute \(w(x)\in T\) and
\(\sigma(x)\in\{-1,0,1\}\) with one of the following certificates:
\[
 \begin{array}{c|c}
 \sigma(x)=1&w(x)x\in T\\
 \sigma(x)=-1&-w(x)x\in T\\
 \sigma(x)=0&w(x)x=0.
 \end{array}
 \tag{2.4}\label{eq:signcertificate}
\]
When the coefficients have the same weak sign, take \(w=1\) and
that sign, with the zero pair assigned sign zero. When \(a,b\) have
opposite nonzero signs, take \(w=(|a|,|b|)\). Direct multiplication gives
\[
 wx=\bigl(\sgn(a)(a^2-b^2d),\,0\bigr).
 \tag{2.5}\label{eq:conecertificate}
\]
The sign of this scalar is exactly \eqref{eq:sign}. Thus the
certificates always exist by a fixed finite calculation.

Certificates of different signs cannot coexist. For example, if
\(ux\in T\) and \(vx=0\), with \(u,v\in T\), then \(v(ux)\in T\)
but equals zero, a contradiction. If \(ux\in T\) and \(-vx\in T\),
then
\[
 v(ux)+u(-vx)=0
\]
is a sum of two elements of \(T\), again impossible. The other
conflicts follow by negation. Consequently any certificate of the
forms \eqref{eq:signcertificate} has the computed sign; there is no
search over possible multipliers.

Write \(P=\{x:\sigma(x)=1\}\) and \(Z=\{x:\sigma(x)=0\}\).
If \(x,y\in P\), with witnesses \(u,v\), then
\[
 uv(x+y)=v(ux)+u(vy)\in T,\qquad
 uvxy=(ux)(vy)\in T.
\]
Thus \(P\) is closed under addition and multiplication.
If \(x,y\in Z\), then \(uv(x+y)=0\); if \(x\in Z\), then
\(u(xz)=z(ux)=0\) for every \(z\). Negation preserves \(Z\).
Hence \(Z\) is an ideal. Finally, if \(x\in P\) and \(z\in Z\),
\[
 uv(x+z)=v(ux)\in T,
\]
so \(P+Z\subseteq P\). Negation reverses nonzero signs.

These explicit certificates prove that sign is well defined modulo
\(Z\), that its trichotomy is total, and that the order is compatible
with addition and multiplication. Nonzero factors have nonzero product,
so the ordered quotient is an integral domain. Scalars \((a,0)\)
have exactly the sign of \(a\), proving that the embedding of \(K\)
is injective and order preserving. The class of \(r\) is positive
and \(r^2=d\). If \(d=e^2\) for some \(e\ge0\) already in \(K\),
the same sign test gives \(\sigma(r-e)=0\), selecting the positive
branch even in a redundant extension.
\end{proof}

Nonnegative square roots are unique in this order: if \(x,y\ge0\)
and \(x^2=y^2\), then \((x-y)(x+y)=0\). Either \(x+y=0\), forcing
both to be zero, or the integral-domain law gives \(x=y\).

In the numerical implementation no set-theoretic quotient is formed:
retain the pairs and use \(\sigma(x-y)=0\) as equality.
The lemma supplies its congruence and order laws.
Starting from \(\Q\), iterate this construction along the finite
radical tower. Every multiplier, norm comparison, and closure
certificate is given by the displayed formulas.

If a field implementation is wanted, it too is explicit. For a
nonzero class \(x=a+br\), put \(n=a^2-b^2d\). If \(n\ne0\),
the conjugate divided by \(n\) is an inverse. If \(n=0\),
nonzeroness forces \(a,b\) to be nonzero with the same sign.
Then \(\sigma(x-2a)=0\) by the opposite-sign case, so the scalar
\(1/(2a)\) is an inverse of \(x\) in the quotient. This handles the
case in which blind conjugate inversion would fail. The main
inequality needs only the ordered-domain laws and positive
denominator pairs.

\section{Dyadic regular polygons and the circle constant}\label{sec:roots}
Define
\[
 c_1=0,\quad s_1=1,\quad \eta_1=i,
\]
and, for \(r\ge1\),
\[
 c_{r+1}=\sqrt{\frac{1+c_r}{2}},\qquad
 s_{r+1}=\sqrt{\frac{1-c_r}{2}},\qquad
 \eta_{r+1}=c_{r+1}+is_{r+1}.
 \tag{3.1}\label{eq:halfangle}
\]
Direct calculation, with positive roots, gives
\[
 c_r^2+s_r^2=1,\qquad
 \eta_{r+1}^2=\eta_r,\qquad
 \eta_r^{\,2^r}=-1.
 \tag{3.2}\label{eq:rootfacts}
\]
For example \(2c_{r+1}s_{r+1}=s_r\), because both sides are
nonnegative and their squares agree. For \(r\ge2\),
\(0<c_r<1\), \(s_r>0\), and \(c_r\ge1/2\).
For \(N=2^r\), \(\eta_r\) is the half-step root and
\(\zeta_r=\eta_r^2\) is the step of the regular \(N\)-gon.
The usual notation would be \(c_r=\cos(\pi/N)\) and
\(s_r=\sin(\pi/N)\), but trigonometric analysis is unnecessary.

\begin{lemma}[Order on the upper half-polygon]\label{lem:semicircle}
For \(N=2^r\), put \(y_j=\Ima(\eta_r^j)\). Then
\[
 y_0=y_N=0,\qquad y_j>0\quad(0<j<N).
 \tag{3.3}\label{eq:upper}
\]
\end{lemma}
\begin{proof}
For \(r=1\), the list is \(1,i,-1\).
The even powers at the next stage are the old powers. Moreover,
\[
 \eta_{r+1}^{\,2j+1}
   =\frac{\eta_r^j+\eta_r^{j+1}}{2c_{r+1}}.
 \tag{3.4}\label{eq:midpoints}
\]
To verify this, use
\(1+\eta_r=2c_{r+1}\eta_{r+1}\) and
\(\eta_{r+1}^2=\eta_r\). The imaginary part of an odd power
is a positive multiple of the sum of two adjacent old imaginary
parts. They are nonnegative, and for each interior odd power at least
one is positive. This proves the finite sign table by induction.
\end{proof}

\subsection{Explicit lower and upper polygonal constants}
For \(r\ge2\), \(N=2^r\), set
\[
 p_r=Ns_r,\qquad q_r=\frac{Ns_r}{c_r}.
 \tag{3.5}\label{eq:pq}
\]
These are half the perimeters of the inscribed and circumscribed
regular \(N\)-gons of the unit circle. The half-angle identities give
\[
 p_{r+1}=\frac{p_r}{c_{r+1}},\qquad
 q_{r+1}=\frac{2p_r}{1+c_r}.
 \tag{3.6}\label{eq:pqrec}
\]
Consequently \(p_r\) increases, \(q_r\) decreases, and
\[
 2\le p_r\le q_r\le q_2=4.
 \tag{3.7}\label{eq:pqbound}
\]
For the upper monotonicity, use
\(2c_r\le1+c_r\); for the lower monotonicity use \(c_{r+1}\le1\).
The bounds at the next stage stay nested, so every later lower bound
lies below every earlier upper bound.

Since \(s_r=p_r/N\le4/N\),
\[
 \begin{split}
 0\le q_r-p_r
 &=\frac{p_r(1-c_r)}{c_r}
  =\frac{p_rs_r^2}{c_r(1+c_r)}
 \le \frac{128}{N^2}.
 \end{split}
 \tag{3.8}\label{eq:pqgap}
\]
Thus the nested intervals \([p_r,q_r]\) have an explicit Cauchy
modulus. Their common named value is denoted by \(\pi\).
This phrase is an abbreviation for the interval inequalities; it does
not use a completeness axiom inside arithmetic.

\subsection{Why this is the ordinary \texorpdfstring{\(\pi\)}{pi}}
Let \(U_N\) be the polygonal interpolation of
\(1,\zeta_r,\ldots,\zeta_r^{N-1},1\), on the uniform parameter grid.
Its speed on each segment is \(2Ns_r=2p_r\le8\).
Normalized midpoint insertion gives
\[
 \norm{U_{2N}-U_N}_\infty=1-c_r
 \le s_r^2\le16/N^2.
 \tag{3.9}\label{eq:circlemesh}
\]
The dyadic tail is at most
\(16N^{-2}\sum_{j\ge0}4^{-j}<32/N^2\), where the geometric
tail has a finite rational bound at every requested precision.
This defines a primitive recursive uniform circle name \(U\).
Its norm is one, its dyadic vertices occur in circular order by
Lemma~\ref{lem:semicircle}, and their mesh tends explicitly to zero.
Equivalently, the polygonal sets converge in Hausdorff distance to
the unit circle.

The inscribed perimeter is \(2p_r\). The limiting map is
\(2\pi\)-Lipschitz, since every \(U_N\) is, while its inscribed
perimeters approach \(2\pi\). Its length is therefore \(2\pi\).
Its inscribed area is \(p_rc_r\), and the circumscribed polygon
has area \(q_r\). Their difference is at most \(192/N^2\),
by \eqref{eq:pqgap} and \(p_r(1-c_r)\le64/N^2\).
They bound the unit disk from inside and outside, so its area is
\(\pi\). This identifies the named value with the usual circle constant
without using isoperimetry.

\section{A rational \texorpdfstring{\(\pi\)}{pi}-name with a primitive recursive modulus}
\label{sec:pi}
For \(k\ge0\), set
\[
 r=k+4,\qquad t=4k+16,\qquad h=2^{-t}.
\]
Starting from \(\widehat c_1=0\), perform exactly \(r-1\) steps
\[
 \widehat c_{j+1}
  =\sd\left(\frac{1+\widehat c_j}{2},t\right).
 \tag{4.1}\label{eq:chat}
\]
Then define
\[
 \bar c=\min(1,\widehat c_r+2h),\qquad
 \boxed{\pi_k=2^r\,\sd(1-\bar c^2,t).}
 \tag{4.2}\label{eq:pilower}
\]
All entries here are rational. Every square-root search is bounded,
so this is a fixed primitive recursive function.

\begin{lemma}\label{lem:pilower}
For every \(k\ge0\),
\[
 \pi_k\le p_{k+4},\qquad
 0\le\pi-\pi_k<2^{-k}.
 \tag{4.3}\label{eq:pilowerbounds}
\]
\end{lemma}
\begin{proof}
For \(0\le y\le x\le1\), the function
\(f(x)=\sqrt{(1+x)/2}\) satisfies the algebraic estimate
\[
 0\le f(x)-f(y)
 =\frac{x-y}{2(f(x)+f(y))}\le\frac{x-y}{2}.
 \tag{4.4}\label{eq:rootcontract}
\]
The denominator estimate follows from
\(f(x),f(y)\ge\sqrt{1/2}\).
Combining \eqref{eq:sdbracket} and \eqref{eq:rootcontract}, induction gives
\[
 0\le c_j-\widehat c_j\le2h.
\]
Hence \(c_r\le\bar c\le1\), \(\bar c-c_r\le2h\), and
\[
 0\le(1-c_r^2)-(1-\bar c^2)\le4h.
\]
For \(a\ge b\ge0\),
\(\sqrt a-\sqrt b\le\sqrt{a-b}\), as is seen by squaring.
Therefore
\[
 0\le p_r-\pi_k
 \le2^r(2\sqrt h+h)
 =2^{-k-3}+2^{-3k-12}.
 \tag{4.5}\label{eq:raderror}
\]
Together with \eqref{eq:pqgap},
\[
 0\le\pi-\pi_k
 \le 2^{-2k-1}+2^{-k-3}+2^{-3k-12}<2^{-k}.
 \tag{4.6}\label{eq:pierror}
\]
The final inequality already holds at \(k=0\), and each ratio to
\(2^{-k}\) is nonincreasing thereafter. Every estimate is an
inequality between finitely coded radicals or rationals.
\end{proof}

The sequence \(\pi_k\) need not be proved monotone. We only use
\[
 \pi_k\le p_{k+4}\le p_R\qquad(R\ge k+4).
 \tag{4.7}\label{eq:pimonouse}
\]

\section{The finite Fourier area--energy identity}\label{sec:fourier}
Fix \(N=2^R\ge4\), and abbreviate
\[
 \eta=\eta_R=c+is,\qquad \zeta=\eta^2,\qquad
 x_k=\Rea(\eta^k),\quad y_k=\Ima(\eta^k).
\]
Let \(z_0,\ldots,z_{N-1}\) be any complex coordinates of a closed
polygon. Real and imaginary parts may be rational or belong to the
finite radical tower already constructed. Set
\[
 E=\sum_{j<N}|z_{j+1}-z_j|^2,\qquad
 A=\frac12\sum_{j<N}\Ima(\overline{z_j}z_{j+1}),
 \tag{5.1}\label{eq:EA}
\]
where indices are cyclic, and define the finite Fourier coefficients
\[
 a_k=\frac1N\sum_{j<N}z_j\zeta^{-jk}.
 \tag{5.2}\label{eq:DFT}
\]
These are finite sums. The only scalar denominator is the positive
integer \(N\).

\begin{lemma}[Finite orthogonality]\label{lem:ortho}
\[
 \sum_{j<N}\zeta^{hj}=
 \begin{cases}
 N,&N\mid h,\\
 0,&N\nmid h.
 \end{cases}
 \tag{5.3}\label{eq:ortho}
\]
\end{lemma}
\begin{proof}
Reduce \(h\) modulo \(N\). If \(0<h<N\), write \(h=2^vb\)
with \(b\) odd and \(v<R\), using a bounded scan of the binary digits.
Put \(M=N/2^{v+1}\). Then
\(\zeta^{hM}=\zeta^{bN/2}=-1\).
In every consecutive block of \(2M\) indices, the second half cancels
the first. This is a finite pairing argument; no algebraic inverse
or analytic geometric-series limit occurs.
\end{proof}

Expansion and \eqref{eq:ortho} yield the inverse formula
\(z_j=\sum_{k<N}a_k\zeta^{jk}\). Applying the same orthogonality after
expanding the quadratic expressions gives
\[
 \begin{aligned}
 E&=N\sum_{k<N}|a_k|^2|\zeta^k-1|^2
   =4N\sum_{k<N}|a_k|^2y_k^2,\\
 A&=\frac N2\sum_{k<N}|a_k|^2\Ima(\zeta^k)
   =N\sum_{k<N}|a_k|^2x_ky_k.
 \end{aligned}
 \tag{5.4}\label{eq:parsevalfinite}
\]
For clarity, all off-diagonal terms contain
\(\sum_{j<N}\zeta^{(k-\ell)j}\) and vanish; the diagonal terms
give exactly the displayed factors of \(N\).

\begin{theorem}[Exact nonnegative decomposition]\label{thm:spectral}
\[
 \boxed{cE-4sA
   =4N\sum_{k=1}^{N-1}|a_k|^2y_ky_{k-1}\ge0.}
 \tag{5.5}\label{eq:positiveidentity}
\]
Consequently
\[
 4Ns\,|A|\le NE.
 \tag{5.6}\label{eq:energy}
\]
\end{theorem}
\begin{proof}
The identity
\[
 cy_k-sx_k=\Ima(\eta^{k-1})=y_{k-1}
\]
turns the difference in \eqref{eq:parsevalfinite} into
\eqref{eq:positiveidentity}. The \(k=0\) term is zero.
For \(1\le k<N\), both \(y_k\) and \(y_{k-1}\) are nonnegative
by Lemma~\ref{lem:semicircle}. Thus \(cE\ge4sA\).
Since \(0<c\le1\), \(E\ge cE\ge4sA\).
Conjugating all \(z_j\) negates area and preserves energy, so the
same estimate holds with \(-A\). This gives \eqref{eq:energy}.
\end{proof}

The stronger conclusion \(4N(s/c)|A|\le NE\) is also available.
For the finite arithmetic theorem, \eqref{eq:energy} is sufficient
and avoids an algebraic denominator. Finite Fourier methods in
polygonal isoperimetry have classical antecedents and modern
extensions \cite{kwong}; no novelty claim for the Fourier method
itself is intended here.

\section{An explicit subdivision proves the arithmetic theorem}
\label{sec:subdivision}
Return to the rational polygon \(P\) and upper lengths \(u_i\).
If \(S=0\), each \(u_i=0\), hence each edge is zero and
\(\Area(P)=0\). Thus assume \(S>0\).

Choose an integer
\[
 B\ge\max(1,S)
\]
and a power of two \(N=2^R\) such that
\[
 R\ge k+4,\qquad N\ge m(1+2^kB^2).
 \tag{6.1}\label{eq:meshthreshold}
\]
The least adequate exponent is found by a bounded binary-length
calculation. An explicit, very wasteful alternative is
\(R=k+4+m(1+2^kB^2)\); thus no unbounded minimization is hidden.
Put \(T=N-m>0\), and for each edge set
\[
 k_i=1+\left\lfloor\frac{Tu_i}{S}\right\rfloor.
 \tag{6.2}\label{eq:subcount}
\]
Then \(k_i\ge1\), \(k_iS\ge Tu_i\), and
\[
 \sum_{i<m}k_i
 \le m+\sum_{i<m}\frac{Tu_i}{S}=m+T=N.
 \tag{6.3}\label{eq:subbudget}
\]
Subdivide edge \(i\) into \(k_i\) equal segments. Append repeated
copies of a vertex until the polygon has exactly \(N\) vertices.
Repeated vertices contribute zero edges and zero area. All new
coordinates are rational.

\begin{lemma}[Area and energy under subdivision]\label{lem:refinement}
The new polygon has the same signed area as \(P\), and its energy obeys
\[
 E=\sum_{i<m}\frac{\norm{e_i}^2}{k_i}
 \le\frac{S^2}{N-m}.
 \tag{6.4}\label{eq:subenergy}
\]
\end{lemma}
\begin{proof}
The vertices on edge \(i\) are \(v_i+je_i/k_i\).
Bilinearity and \(\det(e_i,e_i)=0\) show that every successive
determinant is \(\det(v_i,e_i)/k_i\). Summing \(k_i\) of them
gives the original determinant. Thus area is unchanged.

The squared lengths of these subedges sum to
\(\norm{e_i}^2/k_i\). By \eqref{eq:input} and \(k_iS\ge Tu_i\),
\[
 \frac{\norm{e_i}^2}{k_i}
 \le\frac{u_i^2}{k_i}
 \le\frac{S}{T}u_i.
\]
This includes \(u_i=0\), without division by \(u_i\).
Summing proves \eqref{eq:subenergy}.
\end{proof}

Applying \eqref{eq:energy}, with \(p_R=Ns_R\), gives
\[
 \begin{aligned}
 4p_R|\Area(P)|
 &\le NE
 \le\frac{N}{N-m}S^2\\
 &=S^2+\frac{mS^2}{N-m}
 \le S^2+2^{-k}.
 \end{aligned}
 \tag{6.5}\label{eq:finitechain}
\]
The last step uses \eqref{eq:meshthreshold}:
\(N-m\ge m2^kB^2\) and \(S\le B\).
Finally \(\pi_k\le p_{k+4}\le p_R\), so
\eqref{eq:finitechain} proves Theorem~\ref{thm:finite}.

This is the entire core proof. There is no search for a maximizing
polygon, no claim that a maximizer exists, and no passage to an
infinite Fourier expansion.

\section{The sharp polygon inequality as precision arithmetic}
\label{sec:polygonlimit}
Write
\[
 L(P)=\sum_{i<m}\sqrt{\norm{e_i}^2}
\]
for its explicitly computable length name. Let
\(B\ge\max(1,L(P))\) be a rational or integer bound obtained, for
example, by summing the coordinatewise absolute edge increments.
Translation of the first vertex to zero shows that every vertex is
within \(L(P)/2\) of zero: either of the two polygonal routes to
it bounds the distance. Therefore
\[
 |\Area(P)|\le L(P)^2/4\le B^2/4.
 \tag{7.1}\label{eq:crudear}
\]
The determinant estimate used here,
\(|\det(v,w)|\le\norm v\norm w\), follows by expanding
\((v\cdot w)^2+\det(v,w)^2=\norm v^2\norm w^2\).

Choose rational upper lengths by \eqref{eq:sdbracket}, with errors
at most \(2^{-t}\). Then
\[
 L(P)\le S_t\le L(P)+m2^{-t}.
\]
Put \(d=m2^{-t}\). Theorem~\ref{thm:finite} and
Lemma~\ref{lem:pilower} give the explicit estimate
\[
 4\pi|\Area(P)|
 \le L(P)^2+2Bd+d^2+(1+B^2)2^{-k}.
 \tag{7.2}\label{eq:polygonerror}
\]
For any rational \(0<\eps\le1\), choose \(t,k\) so that
\[
 d\le\frac{\eps}{4(B+1)},\qquad
 2^{-k}\le\frac{\eps}{2(1+B^2)}.
 \tag{7.3}\label{eq:polygonprecision}
\]
Then \(d\le1\), and the last three terms in
\eqref{eq:polygonerror} sum to at most \(\eps\).

\begin{corollary}[Sharp polygon inequality]\label{cor:polygon}
In represented-real order,
\[
 4\pi|\Area(P)|\le L(P)^2.
 \tag{7.4}\label{eq:polygonsharp}
\]
\end{corollary}
Here represented-real order means that the displayed inequality
has been proved with an arbitrarily prescribed rational error, with
the choices \eqref{eq:polygonprecision} supplied explicitly.
The proof does not infer a real inequality by a new completeness
or compactness principle.

\section{Area continuity and represented boundaries}\label{sec:curves}
\subsection{A finite area-continuity estimate}
\begin{lemma}\label{lem:areacontinuity}
Let \(P=(p_i)\) and \(Q=(q_i)\) be closed polygons expressed on a
common parameter subdivision, and suppose
\(\norm{p_i-q_i}\le\delta\) at every common vertex. Then
\[
 |\Area(P)-\Area(Q)|
 \le\frac{\delta}{2}\bigl(L(P)+L(Q)\bigr).
 \tag{8.1}\label{eq:areacontinuity}
\]
\end{lemma}
\begin{proof}
Put \(d_i=p_i-q_i\), \(\Delta p_i=p_{i+1}-p_i\), and similarly
for \(q\). Expanding, and shifting the index of a finite cyclic sum,
gives
\[
 2(\Area(P)-\Area(Q))
 =\sum_i\det(d_i,\Delta p_i)
   +\sum_i\det(d_{i+1},\Delta q_i).
 \tag{8.2}\label{eq:areaidentity}
\]
Use the determinant bound term by term, followed by the triangle
inequality. Common subdivision does not change either area or length.
\end{proof}

In particular, a uniformly Cauchy family of polygonal maps with
perimeters bounded by \(B\) has an explicit signed-area Cauchy
modulus: uniform error \(\delta\) gives area error at most \(B\delta\).
If an exact \(m\)-vertex polygon is approximated at its vertices
within \(\eta\), its perimeter changes by at most \(2m\eta\), and
its area changes by at most
\[
 B\eta+m\eta^2
 \tag{8.3}\label{eq:vertexerror}
\]
when the exact perimeter is at most \(B\).
Thus the finite radical or Cauchy-real coordinates used in geometric
explanations can always be replaced by rational coordinates at an
explicitly chosen precision.

\subsection{Polygonal boundary and domain names}
Suppose \(P_n\) is a fixed primitive recursive family of rational
closed polygons, with supplied moduli ensuring
\[
 |\Area(P_n)|\longrightarrow A,\qquad L(P_n)\longrightarrow L.
 \tag{8.4}\label{eq:APnames}
\]
The limits here are names and precision conditions, not new objects
whose existence arithmetic must prove. Let
\(A,L\ge0\), and suppose \(L(P_n),L\le B\).
For a requested error \(\eps\), use the moduli to arrange
\[
 \bigl||\Area(P_n)|-A\bigr|\le\delta,\qquad
 |L(P_n)-L|\le\delta.
\]
Then Corollary~\ref{cor:polygon} gives
\[
 4\pi A\le L^2+(16+2B)\delta
 \tag{8.5}\label{eq:curvetransfer}
\]
as a precision inequality. Choosing
\(\delta\le\eps/(16+2B)\), and allocating an additional arbitrarily
small error to \eqref{eq:polygonsharp} if working entirely with
rational representatives, proves
\[
 \boxed{4\pi A\le L^2.}
 \tag{8.6}\label{eq:curveiso}
\]

For a Jordan-domain name, the polygons are simple domain boundaries
and the presentation specifies their area and perimeter convergence.
Equation \eqref{eq:curveiso} is therefore the ordinary enclosed-area
inequality for that domain. The finite polygon proof never needed
to decide membership of an arbitrary point in an arbitrary real curve.

For this last geometric reading, the presentation must also certify
that \(L\) is the \emph{actual perimeter of the limiting boundary}.
One precise certificate uses the variation lower cut of a boundary
parametrization: every finite inscribed length is at most \(L\), and
specified inscribed lengths approach \(L\) with the supplied modulus.
Hausdorff convergence and convergence of approximating polygon
perimeters alone do not certify this: increasingly fine corrugations
can retain excess length. Without length correctness,
\eqref{eq:curvetransfer} still proves an inequality for the numerical
polygon limits, and equality still forces a circular boundary, but
the converse need not hold for an artificially enlarged length.

To add the rigidity conclusion, also supply a Hausdorff convergence
modulus from the boundaries of \(P_n\) to the represented boundary
\(\Gamma\), and a positive separation of two boundary points in the
nondegenerate case. Section~\ref{sec:shape} proves that equality forces
a circle; with the length-correctness condition just specified, circles
also attain equality.

\subsection{A direct Lipschitz-curve version}\label{sec:lipschitz}
Suppose \(\gamma:[0,1]\to\C\) is closed and \(L\)-Lipschitz,
with a primitive recursive sampling presentation. Let \(P_N\)
interpolate its values at \(j/N\), \(N\) dyadic.
For \(t=(j+u)/N\), \(0\le u\le1\), convex interpolation gives
\[
 \begin{aligned}
 |\gamma(t)-P_N(t)|
 &\le(1-u)\frac{Lu}{N}
       +u\frac{L(1-u)}{N}\\
 &\le \frac{L}{2N}.
 \end{aligned}
 \tag{8.7}\label{eq:lipschitzinterp}
\]
Also \(L(P_N)\le L\). Comparing coarse and refined interpolants
in Lemma~\ref{lem:areacontinuity} supplies the signed-area name \(A\)
and the explicit bound
\[
 |A-\Area(P_N)|\le\frac{L^2}{2N}.
 \tag{8.8}\label{eq:lipschitzarea}
\]
The finite energy obeys \(NE_N\le L^2\).
Thus \eqref{eq:energy} and \(p_R\to\pi\) directly prove
\(4\pi|A|\le L^2\), with explicit precision choices.
If \(L\) is the actual perimeter and the supplied parametrization is
normalized arclength, this is the usual curve statement.

\section{Perimeter as a lower real: no length oracle}\label{sec:lower}
One can avoid a general perimeter convergence modulus. Suppose a
fixed primitive recursive curve presentation supplies rational samples,
a uniform continuity modulus, closure, and an integer variation bound
\(B\ge1\). The bound means that every finite inscribed polygon has
perimeter at most \(B\).

The continuity modulus gives polygonal interpolants with explicit
uniform error, and Lemma~\ref{lem:areacontinuity} gives a signed-area
name \(A\). Define the \emph{lower-real perimeter} by the rational
lower bounds of finite inscribed polygon lengths. Its square has the
analogous lower cut.

The statement \(4\pi|A|\le L^2\) means:
given a rational \(q<4\pi|A|\) and a finite positive-margin certificate,
construct an inscribed polygon and a rational lower bound on its
squared perimeter that is greater than \(q\).

Here is a bounded witness construction. A certificate supplies a
rational \(g>0\) with
\[
 q+g\le4\pi|A|.
 \tag{9.1}\label{eq:gapcert}
\]
Concretely, if \(a_r\) is a rational area approximation with certified
error \(\eps_r\), a rational comparison
\[
 q+g\le4\pi_j\max(0,|a_r|-\eps_r)
\]
is such a finite certificate. The explicit approximation moduli make
these certificates adequate for every strict rational lower bound.
If \(q<0\), any polygon suffices. Otherwise choose the area
approximation so that
\[
 |A-\Area(P_n)|\le g/64.
\]
Since \(4\pi\le16\),
\[
 4\pi|\Area(P_n)|\ge q+3g/4.
 \tag{9.2}\label{eq:marginremain}
\]
The exact inscribed vertices may be Cauchy names. Write \(P=P_n\),
let \(m\) be its number of vertices, and approximate them by a rational
polygon \(Q\) within \(\eta\), chosen so that, for \(d=2m\eta\),
\[
 d\le1,\qquad
 16(B\eta+m\eta^2)\le g/16,\qquad
 2(B+1)d\le g/16.
 \tag{9.3}\label{eq:lengthinterval}
\]
These are explicit rational thresholds; decreasing \(\eta\) to a
dyadic number satisfying them is bounded. Then
\(L(Q)\le B+1\), and \eqref{eq:vertexerror} gives
\[
 4\pi|\Area(Q)|\ge q+11g/16.
\]
Choose \(k\) with
\((1+(B+1)^2)2^{-k}\le g/16\).
Compute rational edge-length intervals for \(Q\) of total width \(e\),
with \(e\le1\) and \(e(2B+3)\le g/16\), and let their totals be
\(R_Q\le L(Q)\le S_Q\). Then
\(S_Q^2-R_Q^2\le g/16\). The finite theorem gives
\[
 \begin{aligned}
 S_Q^2
 &\ge4\pi_k|\Area(Q)|-2^{-k}\\
 &\ge4\pi|\Area(Q)|-(1+(B+1)^2)2^{-k}
 \ge q+10g/16.
 \end{aligned}
\]
Hence \(R_Q^2\ge q+9g/16\).

The approximating polygon \(Q\) need not itself be inscribed.
To produce a valid member of the curve's length lower cut, output
\[
 R=\max(0,R_Q-d).
\]
The perimeter perturbation bound proves \(R\le L(P)\), and
\[
 R_Q^2-R^2\le2(B+1)d\le g/16.
\]
Thus \(R^2\ge q+g/2>q\), now for a genuinely inscribed polygon.
Every index and precision has been bounded by the supplied continuity
modulus, \(B\), the vertex count, and the rational gap \(g\).

Dyadic partitions suffice to generate the usual variation lower cut.
For a given finite partition, round its parameters to a sufficiently
fine dyadic grid using the uniform continuity modulus. The corresponding
finite length sum changes by at most twice the number of sample points
times their pointwise error. Passing to the full dyadic refinement
only increases length. This is a finite triangle-inequality argument.

Thus the inequality itself does not require computability of the
length or extraction of an arclength parametrization.

\section{An explicit stability estimate and circle rigidity}
\label{sec:rigidity}
\subsection{The two surviving Fourier modes}
Retain the notation of Section~\ref{sec:fourier}, and put
\[
 q_R=\frac{Ns}{c},\qquad
 D=NE-4q_RA.
\]
Multiplying \eqref{eq:positiveidentity} by \(N/c\) yields
\[
 D=\sum_{k=2}^{N-1}\mu_k|a_k|^2,\qquad
 \mu_k=\frac{4N^2}{c}\,y_ky_{k-1}>0.
 \tag{10.1}\label{eq:stabilityspectral}
\]
The modes \(k=0,1\) vanish. For the remainder
\(R_j=z_j-a_0-a_1\zeta^j\), finite weighted Cauchy--Schwarz gives
\[
 |R_j|^2\le D\sum_{k=2}^{N-1}\frac1{\mu_k}.
 \tag{10.2}\label{eq:weightedcauchy}
\]
This is a finite algebraic inequality: apply the usual finite
sum-of-squares proof of Cauchy--Schwarz to
\(\sqrt{\mu_k}a_k\zeta^{jk}\) and \(1/\sqrt{\mu_k}\).
All denominators are positive by Lemma~\ref{lem:semicircle}.

The reciprocal sum telescopes exactly. Namely
\[
 \frac{x_{k-1}}{y_{k-1}}-\frac{x_k}{y_k}
 =\frac{s}{y_ky_{k-1}},
\]
because the numerator is the imaginary part of \(\eta\).
The endpoints are \(x_1/y_1=c/s\) and
\(x_{N-1}/y_{N-1}=-c/s\). Hence
\[
 \sum_{k=2}^{N-1}\frac1{\mu_k}
 =\frac{c^2}{2N^2s^2}\le\frac18,
 \tag{10.3}\label{eq:reciprocal}
\]
using \(Ns=p_R\ge2\). We have proved the finite quantitative statement
\[
 \boxed{\max_j|z_j-a_0-a_1\zeta^j|^2\le D/8.}
 \tag{10.4}\label{eq:stability}
\]
The maximum here is simply a bounded list assertion. This estimate is
strong enough to prove rigidity without selecting subsequences.

\subsection{Explicit passage to a Lipschitz curve}
For a closed \(L\)-Lipschitz curve, translate so that \(\gamma(0)=0\).
Closure gives \(|\gamma(t)|\le L\min(t,1-t)\le L/2\), and so
\[
 |\Area(P_N)|\le L^2/4.
 \tag{10.5}\label{eq:areaboundlip}
\]
The slightly improved version of \eqref{eq:pqgap} is
\[
 0\le q_R-\pi\le q_R-p_R\le64/N^2.
 \tag{10.6}\label{eq:gap64}
\]
Indeed \(c^2\ge1/2\), \(c\ge1/2\), hence
\(c(1+c)=c+c^2\ge1\), and \(p_Rs^2\le64/N^2\).

Put \(\Delta=L^2-4\pi A\ge0\), using signed area; reverse
orientation when applying the result to positive geometric area.
Using \(NE\le L^2\), \eqref{eq:lipschitzarea},
\eqref{eq:areaboundlip}, and \eqref{eq:gap64}, we obtain
\[
 \begin{aligned}
 D
 &\le\Delta+4\pi|A-\Area(P_N)|
       +4(q_R-\pi)|\Area(P_N)|\\
 &\le\Delta+\frac{8L^2}{N}+\frac{64L^2}{N^2}
 \le\Delta+\frac{24L^2}{N}.
 \end{aligned}
 \tag{10.7}\label{eq:Derror}
\]
This argument does not decide the sign of the finite area approximation.
Equation \eqref{eq:stability} consequently gives
\[
 \max_j|\gamma(j/N)-b_N-a_NU(j/N)|^2
 \le\Delta/8+3L^2/N,
 \tag{10.8}\label{eq:gridstability}
\]
where
\[
 b_N=\frac1N\sum_{j<N}\gamma(j/N),\qquad
 a_N=\frac1N\sum_{j<N}\gamma(j/N)\overline{U(j/N)}.
 \tag{10.9}\label{eq:means}
\]

For any \(K\)-Lipschitz function \(f\), its successive dyadic left-endpoint
averages differ by at most \(K/(4N)\). This follows by pairing each
old sample with its new midpoint, whose distance is at most \(K/(2N)\).
Summing the dyadic geometric tail gives an error at most \(K/(2N)\).
The functions \(\gamma\) and \(\gamma\overline U\) have Lipschitz
constants \(L\) and \(5L\), respectively, since \(U\) is
\(8\)-Lipschitz and \(|\gamma|\le L/2\).
Thus the means define explicit names \(b,a\) satisfying
\[
 |b-b_N|\le L/(2N),\qquad
 |a-a_N|\le5L/(2N),\qquad |a|\le L/2.
 \tag{10.10}\label{eq:meanerror}
\]

For a represented parameter \(t\), choose a rational \(t'\) within
\(1/(2N)\) of \(t\), and round \(t'\) to the nearest grid point
\(j/N\), clamping at the endpoints. Then \(|t-j/N|\le1/N\).
This uses only rational rounding, not a floor operation on an
arbitrary real. Combine the sample estimate, the two
mean errors, and
\(|\gamma(t)-\gamma(j/N)|\le L/N\),
\(|U(t)-U(j/N)|\le8/N\). The result is
\[
 \boxed{|\gamma(t)-b-aU(t)|
 \le\sqrt{\Delta/8+3L^2/N}+8L/N.}
 \tag{10.11}\label{eq:uniformstability}
\]
The coefficients and all bounds have primitive recursive names.

\begin{theorem}[Effective equality for a Lipschitz presentation]
\label{thm:equalitylip}
If \(L^2=4\pi A\), then
\[
 \gamma(t)=b+aU(t)
\]
in uniform represented equality. Given an integer bound \(B\ge L\)
and a rational \(\eps>0\), it is enough to take a dyadic
\[
 N\ge\max\left(4,\frac{48B^2}{\eps^2},\frac{16B}{\eps}\right)
 \tag{10.12}\label{eq:rigiditymodulus}
\]
to make the right side of \eqref{eq:uniformstability} at most \(\eps\).
\end{theorem}

Indeed the square-root term is at most \(\eps/4\), and the last
term at most \(\eps/2\). A supplied positive separation of two curve
points implies \(|a|>0\), since their distance is at most \(2|a|\).
The boundary is therefore a nondegenerate circle.

\subsection{Equality for general polygonal boundary names}
\label{sec:shape}
The preceding proof does not force us to extract an arclength
parametrization of the limiting curve. Apply it separately to each
rational polygon \(P_n\).

Choose rational edge upper bounds \(u_i>0\) so accurately that
\[
 0\le S_n-L(P_n)\le2^{-n},\qquad S_n=\sum_i u_i.
\]
Traverse edge \(i\) during the rational time interval of length
\(u_i/S_n\). The resulting polygonal map is \(S_n\)-Lipschitz.
Zero edges cause no difficulty. Reverse the traversal when its rational
signed area is negative. Its deficit is
\[
 \Delta_n=S_n^2-4\pi|\Area(P_n)|\ge0.
\]
By \eqref{eq:uniformstability}, it lies, in uniform distance and hence
Hausdorff distance, within \(\sqrt{\Delta_n/8}\) of an explicit
possibly degenerate circle \(C_n\). The precise reading is that for
every additional rational error the finite bound in
\eqref{eq:uniformstability} supplies an adequate sample size.

Now suppose the boundary name satisfies
\[
 \dH(P_n,\Gamma)\to0,\quad
 |\Area(P_n)|\to A,\quad S_n\to L
\]
with supplied moduli, and \(L^2=4\pi A\).
The deficits then tend to zero effectively. For example, if
\(S_n,L\le B\), both numerical errors are at most \(\delta\), then
\(\Delta_n\le(2B+16)\delta\).
Thus we have explicit \(\tau_n\to0\) such that
\[
 \dH(\Gamma,C_n)\le\tau_n.
 \tag{10.13}\label{eq:circlehausdorff}
\]

No subsequence is selected. Compute the coordinate extrema of each
rational polygon by finite minima and maxima. Hausdorff convergence
makes these extrema Cauchy with the same error modulus.
Let \(x_{\min},x_{\max},y_{\min},y_{\max}\) be their names, and define
\[
 b_*=\left(\frac{x_{\max}+x_{\min}}2,
            \frac{y_{\max}+y_{\min}}2\right),\qquad
 \rho=\frac{x_{\max}-x_{\min}}2.
 \tag{10.14}\label{eq:canonicalcircle}
\]
If \(C_n\) has center \(b_n\) and radius \(\rho_n\),
\eqref{eq:circlehausdorff} implies
\[
 |b_*-b_n|\le2\tau_n,\qquad |\rho-\rho_n|\le\tau_n.
\]
Circles whose centers and radii differ by these amounts have
Hausdorff distance at most \(3\tau_n\): pair points with the same
unit-circle parameter. Hence
\[
 \dH\bigl(\Gamma,C(b_*,\rho)\bigr)\le4\tau_n.
 \tag{10.15}\label{eq:finalrigidity}
\]
This proves represented equality of the boundary with that circle.
A positive two-point separation excludes \(\rho=0\).

In the lower-real setting of Section~\ref{sec:lower}, equality supplies
the missing perimeter modulus. It identifies the length cut with the
name \(L=\sqrt{4\pi|A|}\). For an inscribed \(P_n\), write
\(\ell_n=L(P_n)\le L\). The polygon inequality gives
\[
 (L-\ell_n)^2\le L^2-\ell_n^2
 \le4\pi\bigl||A|-|\Area(P_n)|\bigr|
 \le16|A-\Area(P_n)|.
 \tag{10.16}\label{eq:equalitylengthmod}
\]
Thus the area modulus yields a perimeter modulus under equality,
without dividing by \(L\) or deciding whether \(L=0\).
The preceding polygon-by-polygon rigidity argument applies.

\subsection{Circles attain equality}
For a circle of radius \(\rho\), the regular dyadic polygon has
perimeter \(2p_R\rho\) and area \(p_Rc_R\rho^2\).
Section~\ref{sec:roots} therefore gives
\[
 L=2\pi\rho,\qquad A=\pi\rho^2,\qquad L^2=4\pi A.
\]
If any larger constant replaced \(4\pi\), this circle would violate
the inequality; the regular polygon approximations already provide
finite violations once their explicit errors are smaller than the
proposed positive improvement.

\section{Why the human proof stays within \texorpdfstring{\(\ipra\)}{iPRA}}
\label{sec:audit}
The restriction to quantifier-free induction must be audited explicitly.
The following are the finite assertions used in the proof.

\begin{longtable}{@{}p{0.22\textwidth}p{0.70\textwidth}@{}}
\toprule
Construction & Quantifier-free invariant after numerical coding\\
\midrule
\endhead
Rational operations &
The represented denominator is positive; cross-multiplied identities
and inequalities agree with the encoded integer calculations.\\
Square-root scan &
The stored index is accepted and no larger scanned index is accepted.
The latter is a bounded Boolean conjunction.\\
Half-angle table &
Each stored \(c_j,s_j\) satisfies the square identities, the stated
signs, and the displayed rational error bound when approximated.\\
Upper semicircle &
The Boolean conjunction of the sign statements for the finite power
table, propagated by \eqref{eq:midpoints}.\\
Finite orthogonality &
The partial sum over each paired block cancels; multiplication of the
specified powers has the stated exponent.\\
Fourier expansions &
The running finite sums equal their expanded double sums, with the
off-diagonal terms cancelled by the finite orthogonality table.\\
Subdivision &
The accumulated segment count stays within its rational floor bound;
the accumulated area is unchanged and the energy is bounded by the
corresponding partial sum of \(Su_i/T\).\\
Precision selection &
The powers of two satisfy the explicit numerical thresholds in
\eqref{eq:meshthreshold}, \eqref{eq:polygonprecision}, and
\eqref{eq:rigiditymodulus}.\\
Curve comparison &
For the finitely selected samples, the assumed name/modulus conditions
give the required rational error bounds; the remaining calculation is
a finite inequality.\\
\bottomrule
\end{longtable}

Each indexed conjunction or double sum is a primitive recursive fold
over an explicitly bounded list. Induction is applied to equality or
order statements about those fold values, not directly to a formula
quantifying over arbitrary curves or all real numbers. When a
recursion argument needs bounds on several entries at once, store the
bounded table and use its Boolean conjunction as the invariant.
Increasing the coefficient bound or table size is done by an explicit
primitive recursive majorant before the induction.

For completeness, the radical-order induction can be bounded as follows.
Fix the final tower height \(H\), the tower itself, and a bound \(M\ge2\)
on all scalar numerators and denominators in the relevant arrays and
radicands. A height-\(j\) multiplication uses fewer than \(2\cdot5^j\)
rational operations; its recursive splitting has five multiplications
one level lower plus coefficient additions. The sign calculation uses
fewer than \(10\cdot5^j\) rational operations, since it makes at most
three lower sign calls and three lower multiplications.
The fixed finite list of order-law certificate calculations above is
therefore bounded conservatively by
\[
 Q(j)=1000\cdot8^{j+1}.
\]
One unnormalized rational operation changes a scalar magnitude and
denominator bound \(B\) to at most \(2B^2\). Consequently
\[
 F(j,M)=(2M+2)^{\,2^{Q(j)}}
 \tag{11.1}\label{eq:majorant}
\]
is a deliberately generous bound for all coefficients used in those
calculations at height \(j\).

Precompute a backward budget table
\(M_H=M\), \(M_j=F(j+1,M_{j+1})\). Induct upward on \(j\) using the
Boolean conjunction of the sign and order-law instances for all
height-\(j\) arrays with scalar numerators and denominators bounded by
\(M_j\). The explicit next-level certificate invokes only lower-level
instances inside the preceding budget. Array dimension is \(2^j\),
so the enumeration is bounded by an explicit power of a finite
scalar-code bound. The invariant is quantifier free.
This fixes the final height and coefficient bound as parameters;
it does not use induction on the unbounded assertion that every
element of every tower satisfies the ordered-field axioms.
Afterward the arbitrary finite bounds may be universally introduced.

The algebraic sign recursion from Section~\ref{sec:arithmetic} reduces
to comparisons one tower level lower. Lemma~\ref{lem:orderedextension}
proves its order laws by the explicit multipliers
\eqref{eq:signcertificate}--\eqref{eq:conecertificate}; it does not
appeal to an ambient real field. To internalize a finite calculation,
store its operand arrays and these witnesses, including the
lower-level comparisons required by each witness. Their size bounds
are computed by the same fixed recursion. Checking the table in
dependency order is a bounded Boolean invariant.
All elementary order arguments in the spectral proof are thus
statements about finite coefficient codes. Positive
denominators may be cleared, so no order axiom for an unrepresented
real field is added.

There is no induction over a genuinely analytic predicate. In
particular, \emph{taking a limit} means: given a precision input,
compute the finite indices displayed in the proof and obtain the
corresponding rational inequality. Likewise, equality of two represented
curves is the assertion of all such precision inequalities, proved by
introducing the precision variable and performing the finite argument.
Universal introduction is allowed in intuitionistic logic; it is
unrestricted induction that is unavailable.

The curve statement is a schema for each fixed family of primitive
recursive sample functions and moduli, using their defining equations
and the stated name hypotheses. It is not a purported single
primitive recursive evaluation function for all primitive recursive
program codes. An alternative formulation can carry finite computation
traces with an explicitly supplied primitive recursive time bound.

The proof uses no excluded middle for undecidable statements, no
Markov principle, no countable choice, no Bolzano--Weierstrass principle,
and no variational existence theorem. Decisions about signs of the
finite rational or radical expressions are supplied by terminating
primitive recursive comparisons. This list describes the actual proof,
not a claim that every finite-looking argument automatically belongs
to \(\ipra\).

\section{A concrete path to complete Lean formalization}
\label{sec:lean}
\subsection{Repository snapshot and reusable infrastructure}
The inspected ProveIt snapshot is
\begin{center}\small\ttfamily
329161f8da0b319d51ba9f6246c4b19d243f7193.
\end{center}
Its root configuration pins Lean \texttt{v4.32.0}, with Mathlib revision
\begin{center}\small\ttfamily
81a5d257c8e410db227a6665ed08f64fea08e997.
\end{center}
These are source observations, not a claim to have rebuilt the
repository. The source directory of the Lake library
\texttt{ConciseIPRA} is
\begin{center}\small\ttfamily
Logic/PrimitiveRecursiveArithmetic/ConciseIPRA/Lean.
\end{center}

\texttt{ConciseIPRA.PR} defines primitive recursive function symbols
\texttt{PR}, terms \texttt{PTm}, the theory \texttt{PRAAx}, and the
derived rule \texttt{PRA.qfInd}. Its quantifier-free restriction is
expressed by an explicit \texttt{Fm.QF} premise
\cite{proveit-pr}. \texttt{ConciseIPRA.Deduction} defines the embedded
intuitionistic derivability judgment \texttt{Prf Ax Gamma phi}
\cite{proveit-deduction}. Thus the actual completion criterion is
\begin{lstlisting}
Prf PRAAx [] (finiteFormula checkCode)
\end{lstlisting}
for the compiled checker code, not an external theorem whose type
merely mentions Lean's \texttt{Rat} or \texttt{Real}.

The existing theorem \texttt{pra\_iff\_native} can subsequently
transfer the result to the concise native calculus
\cite{proveit-native}. No dependency on the Busy Beaver or
Diophantine-consistency developments is needed.

\subsection{Files supplied in this archive}
\texttt{FiniteData.lean} defines rational points, cyclic edges, signed
area, the bounded square-root approximation, the exact \(\pi_k\)
algorithm, validity of an upper-length list, and the finite target.

\texttt{PrimitiveCoding.lean} specifies concrete natural encodings of
integers, rationals, points, and finite lists, and a total three-input
checker. Malformed data are accepted vacuously, so the target can be
a universal equation. On valid input, acceptance means exactly
\eqref{eq:finite}.

\texttt{InternalGoal.lean} constructs the first-order formula
\[
 \forall p\,\forall u\,\forall k\;
   \operatorname{check}(p,u,k)=0,
 \tag{12.1}\label{eq:internal}
\]
where zero is the success code. The argument
\texttt{checkCode : PR 3} is explicitly required to be the compilation
of that checker. The file does not substitute an arbitrary constantly
successful program. Its only theorem transfers an \emph{already
supplied} internal proof to the native calculus; the premise of this
conditional theorem is the entire unfinished proof obligation.

The files contain no \texttt{sorry}, \texttt{admit}, or added axiom
declarations. They are nevertheless \emph{uncompiled blueprints}:
no Lean or Lake executable was available in this environment.
Absence of placeholders must not be confused with completion of the
main theorem.

\subsection{Implementation sequence and acceptance criteria}
\begin{enumerate}
\item \textbf{Compile the finite codecs and arithmetic.}
Construct fixed \texttt{PR} codes for natural-list decoding, signed
integer arithmetic, rational operations, bounded scans, and Boolean
folds. Prove their equations internally. The mathematical integers
and rationals in the explanatory layer are specifications of these
codes.

\item \textbf{Prove the root approximation library.}
Formalize \eqref{eq:sdbracket}, the half-angle error invariant, and
\eqref{eq:pierror}. This gives a concrete internal correctness theorem
for the code computing \(\pi_k\).

\item \textbf{Formalize the finite positive-radical calculation.}
Implement coefficient-array addition and multiplication and the
sign recursion. Prove the sign and square identities used by
\eqref{eq:halfangle}. Begin with the restricted fixed towers of this
proof; a general real-algebraic-number package is unnecessary.

\item \textbf{Prove the dyadic root table and finite DFT lemmas.}
Establish \eqref{eq:upper}, \eqref{eq:ortho}, and
\eqref{eq:parsevalfinite}. Keep bounded index conjunctions as Boolean
folds so that each invocation of \texttt{PRA.qfInd} has an actual
\texttt{Fm.QF} witness.

\item \textbf{Formalize subdivision and close the finite theorem.}
Prove the rational floor budget, area preservation, and
\eqref{eq:subenergy}; combine with the spectral identity and the mesh
threshold. Produce a derivation of the matrix in \eqref{eq:internal},
then apply universal introduction three times.

\item \textbf{Add precision-indexed geometric corollaries.}
For fixed PR sample families, formalize the finite determinant
identity \eqref{eq:areaidentity}, the error budgets, the lower-real
witness procedure, and the stability telescope. The curve and circle
names should remain first-order formulas with explicit moduli.

\item \textbf{Optionally connect to Mathlib geometry.}
In the external metatheory, prove that the selected effective
Jordan-domain representation agrees with conventional enclosed area
and perimeter. This interpretation layer can use existing analysis
infrastructure, but it must be kept logically distinct from the
internal derivation.
\end{enumerate}

The first meaningful gate is the checked internal finite theorem,
with the checker compilation fixed and its equation audited.
The second is the represented-domain corollary, and the third the
rigidity statement. A separate classical Lean proof, such as the
analytic route described by Samarakkody \cite{samarakkody}, can inform
the interpretation layer; it does not certify the arithmetic strength
of this argument.

\section{Artifact verification and reproducibility}
The archive includes exact rational reference code, an exact radical
Fourier diagnostic, the Lean blueprint, and machine-readable check
reports. The rational checks cover:
\begin{itemize}
\item 25 \(\pi_k\) intervals, compared using an independent rational
Machin-series enclosure;
\item 2,160 rational square-root interval cases;
\item 845 finite polygon instances, including degeneracies, repeated
vertices, orientation reversal, self-intersection, and very large
translations;
\item six expanded subdivisions, checking area preservation and the
energy identity exactly.
\end{itemize}
The radical diagnostic verifies the root and sign tables,
orthogonality, the separate Fourier area and energy formulas, and
\eqref{eq:positiveidentity} for four examples of sizes
\(4,4,8,16\), using exact fractions and recursively represented
radicals. It uses no floating-point mathematics or trigonometric
library.

These are implementation and formula checks, not a replacement for
the universal proof and not an internal \(\ipra\) derivation.
The Machin comparison is an independent diagnostic only; the proof
of the constant in this note is the polygonal construction in
Sections~\ref{sec:roots}--\ref{sec:pi}.

From the archive root, run:
\begin{lstlisting}
python3 code/verify.py
python3 code/dyadic_spectral_check.py
sh build.sh
\end{lstlisting}
The Python scripts require only the standard library.
The build script requires a LaTeX installation with the packages
listed in the source preamble. The compiled PDF is supplied so that
rebuilding is optional.

\begin{thebibliography}{9}
\raggedright
\bibitem{osserman}
R. Osserman, \emph{The isoperimetric inequality},
Bulletin of the American Mathematical Society 84 (1978), 1182--1238.
\href{https://doi.org/10.1090/S0002-9904-1978-14553-4}{DOI}.
\href{https://metaphor.ethz.ch/x/2021/fs/401-3462-00L/sc/extras/Osserman.pdf}{Accessible primary-paper copy}.

\bibitem{kwong}
K.-K. Kwong, \emph{A unified approach to higher order discrete and
smooth isoperimetric inequalities}, Mathematika,
\href{https://doi.org/10.1112/mtk.12238}{doi:10.1112/mtk.12238}.
\href{https://arxiv.org/abs/2111.13352}{arXiv:2111.13352}, version 2,
26 December 2023.

\bibitem{gu-lutz-mayordomo}
X. Gu, J. H. Lutz, and E. Mayordomo,
\emph{Curves That Must Be Retraced}, Sixth International Conference
on Computability and Complexity in Analysis, 2009.
\href{https://drops.dagstuhl.de/entities/document/10.4230/OASIcs.CCA.2009.2267}{Primary publication and paper}.

\bibitem{proveit-pr}
V. Reshetnikov, ProveIt repository,
\texttt{ConciseIPRA/PR.lean}, inspected at commit
\texttt{329161f8da0b319d51ba9f6246c4b19d243f7193}.
\href{https://github.com/VladimirReshetnikov/ProveIt/blob/329161f8da0b319d51ba9f6246c4b19d243f7193/Logic/PrimitiveRecursiveArithmetic/ConciseIPRA/Lean/ConciseIPRA/PR.lean}{Pinned source}.
The root \texttt{lean-toolchain}, \texttt{lakefile.toml}, and
\texttt{lake-manifest.json} at this commit supply the configuration
reported in Section~\ref{sec:lean}.

\bibitem{proveit-deduction}
ProveIt, \texttt{ConciseIPRA/Deduction.lean}, same pinned commit.
\href{https://github.com/VladimirReshetnikov/ProveIt/blob/329161f8da0b319d51ba9f6246c4b19d243f7193/Logic/PrimitiveRecursiveArithmetic/ConciseIPRA/Lean/ConciseIPRA/Deduction.lean}{Pinned intuitionistic deduction source}.

\bibitem{proveit-native}
ProveIt, \texttt{ConciseIPRA/NativeBiInterpretation.lean},
same pinned commit.
\href{https://github.com/VladimirReshetnikov/ProveIt/blob/329161f8da0b319d51ba9f6246c4b19d243f7193/Logic/PrimitiveRecursiveArithmetic/ConciseIPRA/Lean/ConciseIPRA/NativeBiInterpretation.lean}{Pinned native-calculus transfer source}.

\bibitem{samarakkody}
M. Samarakkody, \emph{Formalizing the Classical Isoperimetric
Inequality in the Two-Dimensional Case},
\href{https://arxiv.org/abs/2603.14663}{arXiv:2603.14663},
15 March 2026. Cited as a comparison of formalization approaches;
its proof is not a dependency of this note.
\end{thebibliography}
\end{document}
